\section{Introduction}

\subsection{Motivation}

Software often contains small functions whose logic is commercially or
operationally sensitive: license checks, protocol constants, key-derivation
steps, anti-tamper predicates, and proprietary integer algorithms. The defender
cannot generally prevent a determined local attacker from observing a program:
the attacker has the binary, the CPU, the operating system, and arbitrary
instrumentation. What the defender can do is raise the cost of recovering the
function's logic and make a recovered artifact specific to one build.

Code virtualization is one of the strongest practical transformations available
for this problem. Instead of protecting a native instruction stream in place, the
protector translates the selected function into a private instruction set and
executes that instruction set with an embedded interpreter. The original
instructions disappear from the image, the bytecode can be encrypted and
shuffled, and the interpreter can be diversified per build. The cost is a
substantial runtime overhead and a strict limit on the code that can be lifted.

Commercial virtualizers such as VMProtect, Code Virtualizer, Themida, and
Denuvo are mature, closed products with full instruction-set coverage, packing,
anti-debug, licensing, and multi-platform support. Research systems such as
Tigress, Obfuscator-LLVM, and MOVfuscator explore complementary points in the
design space. Sirus occupies a deliberately narrower point: a reproducible,
source-auditable, zero-dependency implementation for selected x86-64 leaf
functions, built in safe Rust, with a measured evidence trail rather than
vendor claims.

\subsection{Problem statement}

Given an ELF64 or PE32+ AMD64 image $B$ and a selected function $f$ that belongs
to a supported x86-64 subset, Sirus must produce an image $B'$ such that:

\begin{enumerate}
  \item $B'$ preserves the observable integer behavior of $f$ on supported
        inputs and the supported ABI;
  \item the original bytes of $f$ no longer contain a readable native
        implementation;
  \item the replacement is not a single fixed interpreter shared by all builds,
        but a per-build randomized artifact whose opcode assignment, operand
        layout, register encoding, state identifiers, payload encryption, and
        code/data placement depend on a 256-bit seed;
  \item the generated runtime is memory-safe under the Rust type system, uses
        no heap, no libc, and no host memory, and fails closed on malformed
        bytecode;
  \item the result can be reproduced exactly from the input, seed, profile, and
        pinned toolchain, and the evidence for each security and correctness
        claim can be re-executed locally.
\end{enumerate}

The problem is deliberately not "construct an indistinguishability obfuscator."
Barak \etal{} showed that general-purpose obfuscation with strong
virtual-black-box properties is impossible for arbitrary programs
\cite{barak2012}. Practical protectors therefore aim at cost asymmetry,
build-specific diversity, and resistance to specific automated analyses. Sirus
follows that realistic objective.

\subsection{Contributions}

This paper makes the following contributions.

\begin{enumerate}
  \item \textbf{A complete selective virtualizer.} We describe and evaluate an
        end-to-end pipeline that parses and edits ELF64 and PE32+ images, lifts a
        strict x86-64 subset into a bounded IR, flattens control flow, rewrites
        arithmetic, encodes a randomized ISA, synthesizes a freestanding
        \code{no_std} Rust VM, injects it into a new read-execute segment, and
        scrubs the original body.

  \item \textbf{A randomized encrypted ISA.} The virtual instruction set has 48
        semantic operations and four aliases per operation. Each alias has its
        own operand order, one of two register encoding maps, and an immediate
        endianness bit. The runtime decode table is itself encrypted and derived
        from the build seed; the alias layout is not serialized as plaintext in
        the executable.

  \item \textbf{A layered obfuscation stack.} The pipeline combines CFG
        flattening with random state identifiers, opaque guards, decoy blocks,
        bounded linear MBA, a non-linear modular-inverse wrapper, constant
        blinding, a rolling-key ARX body stream, SipHash-based integrity and key
        binding, and masked virtual registers and private memory at rest.

  \item \textbf{A fixed-pass oblivious execution mode.} The optional
        \verb|--constant-time| path executes every manifest block on every pass,
        carries the active path in frontier masks, predicates register writes,
        enqueues both branch successors with complementary masks, and scans the
        whole private stack for loads and stores. We state the exact property it
        provides, the pass bound it assumes, and the microarchitectural effects it
        does not model.

  \item \textbf{An evidence-driven evaluation.} We report 226 tests in debug and
        release modes, Kani and Miri results, more than 40 million recorded
        libFuzzer executions across VM, lifter, ELF, and PE targets, differential
        execution against native code, tamper and forced-return experiments,
        Ghidra and radare2 decompilation, bounded angr and KLEE probes, Creusot
        boundary proofs, three external MBA simplifiers, dudect-bencher,
        Cachegrind, Wine, native Windows, and an MSVC \code{/guard:cf} fixture.
        We also report the negative results: the default branchy profile leaks
        timing, GAMBA and SiMBA reduce parts of the MBA layer, one Ghidra
        decompilation exceeds its budget, and the symbolic probes time out rather
        than prove unsatisfiability.

  \item \textbf{A documented boundary.} We enumerate the supported instruction
        subset, ABI limits, platform limits, unsupported attacks, and the
        assumptions behind the integrity binding. The paper does not claim
        kernel-level anti-debug, trusted signing, a formal Spectre proof, or a
        proof of obfuscation security.
\end{enumerate}

\subsection{Artifact and reproducibility}

The artifact is a Cargo workspace with four shipped crates
(\code{sirus-core}, \code{sirus-vm}, \code{sirus-binfmt},
\code{sirus-protector}), an excluded fuzzing crate, and external verification
harnesses under \path{verification/} and \path{scripts/}. The release is
\code{1.5.9-hardened}. The fast gate is \path{scripts/verify.sh}; the long proof
and fuzzing gate is \path{scripts/prove.sh}. The public Windows CI workflow
executes Wine, native Windows, and MSVC \code{/guard:cf} jobs.

A fixed input, seed, profile, and toolchain produce a bit-identical protected
image. The differential test harness compiles the same C source twice, protects
selected functions with a fixed seed, and compares native and protected results
over thousands of random inputs. The measurement scripts and their outputs are
part of the repository, so a reader can distinguish a reproduced result from a
quoted one.

\subsection{Non-goals}

Sirus is not a whole-program virtualizer, a packer, a licensing system, or a
replacement for a formally audited protector. It does not support external
calls, indirect calls or jumps, writable RIP-relative globals, host pointers,
floating point, x87, AVX-256, other instruction-set extensions, syscalls,
atomics, or arbitrary platform ABIs. It does not implement kernel anti-debug or
guarantee security against an attacker with unlimited local instrumentation. The
default profile is not constant-time. These are design boundaries, not hidden
failure modes: unsupported inputs fail the build, and unsupported runtime states
fail closed with a deterministic status.

\subsection{Organization}

Section~\ref{sec:threat} defines the adversary and objectives.
Section~\ref{sec:background} reviews obfuscation, virtualization, MBA, binary
analysis, and timing measurement. Section~\ref{sec:overview} gives the system
overview. Sections~\ref{sec:frontends}--\ref{sec:runtime} describe the binary
front ends, lifting, obfuscation, and VM runtime.
Section~\ref{sec:constant-time} presents the oblivious mode;
Section~\ref{sec:hardening} describes the optional hardening tier;
Section~\ref{sec:implementation} covers safety engineering.
Section~\ref{sec:evaluation} reports the evaluation, and
Sections~\ref{sec:security}--\ref{sec:conclusion} analyze residual risk and
conclude. The appendices contain the full virtual ISA, the x86-64 subset, the
verification matrix, reproduction commands, and a glossary.
